% =============================================================
% preamble.tex — общая преамбула отчётов по дисциплине
% «Стили и методы программирования. Часть 2»
% Оформление по СТП БГУИР 01-2024; сборка XeTeX (tectonic)
% =============================================================

% Encoding and language
\usepackage{fontspec}
\usepackage[russian]{babel}

% Times New Roman font
\setmainfont{Times New Roman}
\setmonofont{Menlo}[Scale=0.85]

% Page geometry: STP 2.1.1
\usepackage{geometry}
\geometry{left=30mm, right=15mm, top=20mm, bottom=20mm}

% Line spacing: STP 2.1.1
\usepackage{setspace}
\singlespacing

% Allow TeX a bit more flexibility to avoid overflow
\pretolerance=1000
\tolerance=2000
\emergencystretch=3em

% Paragraph indent
\usepackage{indentfirst}
\setlength{\parindent}{1.25cm}

% Page numbers
\usepackage{fancyhdr}
\pagestyle{fancy}
\fancyhf{}
\fancyfoot[R]{\thepage}
\renewcommand{\headrulewidth}{0pt}
\fancypagestyle{plain}{%
  \fancyhf{}%
  \fancyfoot[R]{\thepage}%
  \renewcommand{\headrulewidth}{0pt}%
}

% Math
\usepackage{amsmath,amssymb}

% Graphics
\usepackage{graphicx}
\usepackage{float}
\usepackage{xcolor}

% Tables
\usepackage{booktabs}
\usepackage{tabularx}
\usepackage{multirow}
\usepackage{longtable}

% Captions
\usepackage{caption}
\captionsetup[figure]{
  labelsep=endash,
  justification=centering,
  font={normalsize},
  name=Рисунок
}
\captionsetup[table]{
  labelsep=endash,
  justification=raggedright,
  singlelinecheck=false,
  font={normalsize},
  name=Таблица
}

% Section formatting
\usepackage{titlesec}
% Каждый раздел с новой страницы. \newpage внутри \titleformat срабатывает
% после записи в оглавление и сбивает номера страниц, поэтому \clearpage
% ставится перед самой командой \section.
\usepackage{etoolbox}
\pretocmd{\section}{\clearpage}{}{}
\titleformat{\section}
  {\normalfont\fontsize{14}{18}\selectfont\bfseries}
  {\thesection}{0.5em}{\MakeUppercase}
\titlespacing*{\section}{1.25cm}{0pt}{14pt}

\titleformat{\subsection}
  {\normalfont\fontsize{14}{18}\selectfont\bfseries}
  {\thesubsection}{0.5em}{}
\titlespacing*{\subsection}{1.25cm}{14pt}{14pt}

\titleformat{\subsubsection}
  {\normalfont\fontsize{14}{18}\selectfont\bfseries}
  {\thesubsubsection}{0.5em}{}
\titlespacing*{\subsubsection}{1.25cm}{14pt}{14pt}

% Unnumbered sections
\newcommand{\sectioncenter}[1]{%
  \newpage
  \phantomsection
  \addcontentsline{toc}{section}{#1}
  \begin{center}
    \textbf{\MakeUppercase{#1}}
  \end{center}
  \vspace{14pt}
}

% Table of contents formatting
\usepackage{tocloft}
\renewcommand{\cftsecfont}{\normalfont}
\renewcommand{\cftsecpagefont}{\normalfont}
\renewcommand{\cftsubsecfont}{\normalfont}
\renewcommand{\cftsubsecpagefont}{\normalfont}
\renewcommand{\cftsecleader}{\cftdotfill{\cftdotsep}}
\renewcommand{\cftdotsep}{1}
% babel переопределяет \contentsname при \begin{document}, поэтому через captionsrussian
\addto\captionsrussian{\renewcommand{\contentsname}{СОДЕРЖАНИЕ}}
\renewcommand{\cfttoctitlefont}{\hfill\normalfont\bfseries}
\renewcommand{\cftaftertoctitle}{\hfill}

% Lists
\usepackage{enumitem}
\setlist[itemize]{noitemsep, topsep=0pt, leftmargin=1.25cm, label=--}
\setlist[enumerate]{noitemsep, topsep=0pt, leftmargin=1.25cm}

% Code listings: fvextra корректно работает с Unicode под XeTeX
% (listings переставляет символы «», —, № в комментариях)
\usepackage{fvextra}
\fvset{
  frame=single,
  fontsize=\small,
  numbers=left,
  numbersep=6pt,
  breaklines=true,
}
\newcommand{\code}[1]{\texttt{#1}}

% Hyperlinks
\usepackage{hyperref}
\hypersetup{
  colorlinks=true,
  linkcolor=black,
  urlcolor=blue,
  citecolor=black,
}

% Number figures and tables within sections
\numberwithin{figure}{section}
\numberwithin{table}{section}
\numberwithin{equation}{section}

% Титульный лист: \makelabtitle{номер}{тема}{дата}
\newcommand{\makelabtitle}[3]{%
  \begin{titlepage}
  \thispagestyle{empty}
  \begin{center}

  МИНИСТЕРСТВО ОБРАЗОВАНИЯ РЕСПУБЛИКИ БЕЛАРУСЬ

  \vspace{0.5em}

  Учреждение образования

  \vspace{0.5em}

  \textbf{<<Белорусский государственный университет информатики и радиоэлектроники>>}

  \vspace{2em}

  Факультет компьютерных систем и сетей

  \vspace{0.5em}

  Кафедра программного обеспечения информационных технологий

  \end{center}

  \vspace{3em}

  \begin{center}
  \textbf{ОТЧЕТ}

  \vspace{0.5em}

  по практической работе №#1

  \vspace{0.5em}

  по дисциплине

  \vspace{0.5em}

  <<Стили и методы программирования. Часть 2 (МДН)>>

  \vspace{0.5em}

  на тему

  \vspace{0.5em}

  <<#2>>
  \end{center}

  \vspace{5em}

  \begin{flushright}
  \begin{minipage}{0.5\textwidth}
  \textbf{Выполнил:}\\
  Лебедевич Артем Владимирович\\
  магистрант 2 курса группы 556301

  \vspace{0.6em}

  \textbf{Преподаватель:}\\
  Парамонов Антон Иванович\\
  кандидат технических наук\\
  доцент

  \vspace{0.6em}

  \textbf{Дата:} #3
  \end{minipage}
  \end{flushright}

  \vfill

  \begin{center}
  Минск, 2026~г.
  \end{center}

  \end{titlepage}
}
